




\documentclass[11pt]{article}

\usepackage[a4paper,margin=1in]{geometry}
\usepackage[T1]{fontenc}
\usepackage[utf8]{inputenc}
\usepackage{lmodern}
\usepackage{amsmath,amssymb,amsthm,bm,mathtools}
\usepackage{enumitem}
\usepackage{array}
\usepackage{longtable}
\usepackage{tabularx}
\usepackage{booktabs}
\usepackage{xurl}
\usepackage[expansion=false]{microtype}
\setlength{\emergencystretch}{3em}
\usepackage{xcolor}
\usepackage{hyperref}
\usepackage[nameinlink,noabbrev]{cleveref}

\usepackage[backend=biber,style=numeric,sorting=none,maxnames=6,uniquename=false,url=false]{biblatex}
\AtEveryBibitem{\clearfield{note}\clearfield{annotation}\clearfield{keywords}}
\AtBeginBibliography{\raggedright\small\setlength{\bibitemsep}{0pt}}
\addbibresource{cmotion_foundational_general.bib}
\addbibresource{cmotion_primary_research.bib}
\addbibresource{cmotion_internal_manuscripts.bib}

\newtheorem{theorem}{Theorem}[section]
\newtheorem{proposition}[theorem]{Proposition}
\newtheorem{corollary}[theorem]{Corollary}
\theoremstyle{definition}
\newtheorem{definition}[theorem]{Definition}
\newtheorem{remark}[theorem]{Remark}

\newcommand{\cmotion}{\mbox{C-Motion}}
\newcommand{\ElEmEr}{\mbox{ElEmEr}}
\newcommand{\cmc}{c}
\newcommand{\dd}{\mathrm{d}}
\newcommand{\ii}{\mathrm{i}}
\newcommand{\RR}{\mathbb{R}}
\newcommand{\CC}{\mathbb{C}}
\newcommand{\UU}{\mathrm{U}(1)}
\newcommand{\SU}{\mathrm{SU}(2)}
\newcommand{\SO}{\mathrm{SO}(3)}
\newcommand{\Order}{\mathcal{O}}
\newcommand{\Tr}{\operatorname{Tr}}
\newcommand{\ket}[1]{\left|#1\right\rangle}
\newcommand{\bra}[1]{\left\langle#1\right|}

\newcommand{\statusDerived}{\textcolor{green!50!black}{\textbf{Derived}}}
\newcommand{\statusModelDerived}{\textcolor{cyan!55!black}{\textbf{Model-derived}}}
\newcommand{\statusStructural}{\textcolor{blue!55!black}{\textbf{Structural}}}
\newcommand{\statusAssumed}{\textcolor{orange!80!black}{\textbf{Assumed}}}
\newcommand{\statusConditional}{\textcolor{yellow!60!black}{\textbf{Conditional}}}
\newcommand{\statusNegative}{\textcolor{red!65!black}{\textbf{Negative}}}
\newcommand{\statusConjectural}{\textcolor{magenta!65!black}{\textbf{Conjectural}}}
\newcommand{\statusDeferred}{\textcolor{gray!70!black}{\textbf{Deferred}}}
\newcommand{\statusOpen}{\textcolor{red!75!black}{\textbf{Open}}}
\newcommand{\statusSharpOpen}{\textcolor{red!35!black}{\textbf{Open (sharp)}}}
\newcommand{\statusRetracted}{\textcolor{red!80!black}{\textbf{Retracted}}}

\newcolumntype{L}{>{\raggedright\arraybackslash}X}
\newcolumntype{P}[1]{>{\raggedright\arraybackslash}p{#1}}

\hypersetup{
  colorlinks=true,
  linkcolor=blue!50!black,
  citecolor=blue!50!black,
  urlcolor=blue!50!black,
  pdftitle={Quantum Foundations in C-Motion: From Recurrent Field Organization to Effective Quantum Dynamics},
  pdfauthor={Toma Kavonius}
}

\title{\textbf{Quantum Foundations in \cmotion:\\
From Recurrent Field Organization to Effective Quantum Dynamics}}
\author{Toma Kavonius}
\date{24 November 2026}


\begin{document}
\maketitle

\begin{center}
\small\itshape
Scope and limitations.  This manuscript presents a bounded effective route from the
C-Motion framework toward one-particle quantum dynamics.  Imported normalization,
conditional operator assumptions, and unresolved quantum-measurement completion are
identified explicitly.
\end{center}
\vspace{0.6em}

\begin{abstract}
The C-Motion framework asks whether recurrent organization of a field
propagating locally at speed $c$ can support an effective one-particle quantum
description without taking point particles or complex wavefunctions as primitives.
We examine a restricted free operator class rather than a selected microscopic
ElEmEr law.  Given a translation-invariant Hermitian local first-order sector,
isotropic $\pm c$ characteristics, an unsplit nonzero gap, and an independently
assumed even-squared-symbol condition, its matrices satisfy Clifford relations.
The resulting massive square-root branches admit a controlled low-momentum
Schr\"odinger envelope.  A reconstructed quadratic branch action then yields a
continuous phase charge and translation charges while leaving the absolute
action scale $J_0$ undetermined; neither physical $Q=1$ nor $J_0=\hbar$ is
derived.

Within this free effective sector, Fourier uncertainty and superposition
follow.  Tunnelling and the hydrogen gross spectrum additionally require
supplied interaction and bound-state reductions.  Spinorial $4\pi$ return is
available as a representation structure, not a derived matter identity or
exchange law.  Physical conjugate closure branches offer a conjectural
antimatter interpretation, but complex conjugation alone does not establish
physical charge conjugation or opposite gauge current.  The counterexamples
and exclusions identify why local-$c$ propagation and a gap alone cannot select
the required operator or physical complex/positive-frequency structure.
Microscopic closure dynamics, gap and branch selection, the action quantum,
fermionic exchange, Born probabilities, and definite measurement outcomes
remain unresolved.  Recovering established Dirac and Schr\"odinger mathematics
conditionally is a consistency analysis, not independent empirical validation
of the proposed field ontology.
\end{abstract}

\medskip

\section{Introduction: from recurrent field organization to quantum dynamics}
\label{sec:scope}

Quantum mechanics is extraordinarily successful at predicting amplitudes, spectra, and
measurement statistics.  A deeper physical question remains: what sort of underlying motion
could make its one-particle wave dynamics natural?  \cmotion{} approaches that question from
a field-first picture.  It does not begin by declaring a complex wavefunction, a particle
Hamiltonian, or an action quantum to be primitive.  Instead it asks how much of those
structures can arise once local-$c$ propagation becomes recurrently organized.

The Main Framework supplies the common starting point \cite{CmotionMainFramework2026}.  The
proposed underlying dynamical field/state structure is \textbf{Elemental Motion Resonance
(ElEmEr)}.  Its resolved primitive changes propagate through physical three-space at local
speed $c$.  A primitive succession parameter $s$ records the before--after ordering of that
change; it is an evolution parameter, not a fourth spatial-like dimension and not ordinary
clock time.  A \emph{closure} is a localized or otherwise persistent recurrent organization of
the field, while an open mode propagates without such return.  In this paper a \emph{mode}
means a solution or spectral component of an effective linear evolution, and a \emph{branch}
means one of its separated spectral families, such as the positive- and negative-frequency
families obtained below.

A closure should not be interpreted as necessarily arising from a
single propagation component bending into a closed path. Investigated
kinematic pictures include a localized ElEmEr field packet returning
and potentially overlapping itself, and a coherent recurrent
organization containing oppositely directed propagation contributions
of one underlying field. These are distinct descriptions whose
physical equivalence has not been established. Their directional
momentum contributions may balance while primitive local propagation
remains at $c$, permitting rest-like collective motion with continuing
internal activity.

Such kinematic possibilities do not themselves establish
self-confinement. In particular, cancelling net momentum does not
automatically cancel the stresses relevant to an isolated stationary
closure. The required stress balance, microscopic interaction, capture,
finite-size selection and stability remain unresolved. Neither picture
requires every closure to consist of exactly two waves or permanent
field constituents.

This gives the quantum discussion a concrete starting picture.  A particle-like object is not
treated as a primitive point carrying an attached wave.  It is instead approached as an
organized field state whose primitive propagation can remain local at $c$ while the recurrent
pattern has its own slower collective motion and internal return structure.  Carrier-like phase
variation and a slowly varying envelope are therefore two descriptions of one field
configuration, not two material substances.

Within an effective recurrent mode, phase may be pictured as keeping track of position within a
repeatable cycle.  Only phase differences and phase evolution enter the observable relations
used below; shifting the common phase origin is a symmetry of the reconstructed branch action.
This physical picture motivates the later phase/action construction without claiming that a
unique microscopic phase variable or complex state space has already been selected by the
underlying field dynamics.

It is equally important to keep three physical layers separate: linear propagation within the
effective branch, formation and stabilization of a persistent closure, and interaction with an
environment or apparatus.  The same symbol $\Psi$ can participate in descriptions of more than
one layer, but success in the linear branch does not by itself solve closure formation or
measurement.

The quantum problem can then be stated physically:
\begin{quote}
\emph{If primitive propagation remains local at $c$, can recurrent internal organization
produce a phase/action structure, a rest-like energy scale, spinorial rotation properties,
and a low-energy envelope whose dynamics becomes Schrödinger-like?  If so, where do genuinely
quantum ingredients such as the absolute action scale, action quantization, and measurement
probabilities still have to enter?}
\end{quote}

The constructive chain developed here is
\begin{equation}
\boxed{
\begin{gathered}
\text{local-$c$ field organization}
\longrightarrow
\text{first-order spinorial sector}
\longrightarrow
\text{gapped branch}\\
\longrightarrow
\text{effective mass scale}
\longrightarrow
\text{Schrödinger envelope}\\
\longrightarrow
\text{action/translation charges}
\longrightarrow
\text{one-particle quantum regime.}
\end{gathered}}
\label{eq:readerchain}
\end{equation}
The first arrow is not yet a microscopic derivation: the effective first-order operator and
the nonzero gap remain additional assumptions.  The value of the chain is that, once those
upstream structures are supplied, several later steps are no longer independent quantum
postulates.

The common notions of primitive succession, local-$c$ propagation, operational time,
field-first interaction, and the generic distinction between energy and action are taken from
the Main Framework rather than rederived here \cite{CmotionMainFramework2026}.  Local gauge
architecture and Maxwell dynamics are treated separately in
\emph{Gauge Architecture and Constraints in C-motion I: Recurrent Closure,
Native Constraints, and Conditional Gauge Structure} \cite{CmotionSMStructureI2026}.  Born probabilities, apparatus coupling, and definite
single outcomes remain part of the separate, publication-deferred measurement-completion problem.  The formal scientific-status and provenance
bookkeeping for the results used here is collected in \Cref{app:status}; the main text is
organized around the physical argument.

\subsection{Contribution relative to established theory}
\label{sec:contribution}

The algebraic link between a relativistic first-order multicomponent wave
operator, anticommuting matrices and a massive dispersion relation is part of
the established Dirac/Clifford framework \cite{Dirac1928Electron,LawsonMichelsohn1989SpinGeometry}.
The nonrelativistic reduction of that framework is likewise established
\cite{FoldyWouthuysen1950,Sakurai2017Modern}, as is the correspondence between
continuous action symmetries and conserved charges \cite{Noether1918}.
The present paper does \emph{not} claim originality for those algebraic or
low-energy formulas, nor does their reproduction count as independent
experimental support for \cmotion{}.

The contribution assessed here is more specific.  First, a recurrent-field
interpretation motivates a restricted effective first-order model class,
but the actual operator and its gap are declared as unproved inputs.
Second, \Cref{sec:antecedents,sec:clifford} makes the hypotheses of the
Clifford reduction explicit and exhibits why an independent even-squared-symbol
condition is necessary in the tested class; local speed $c$ and a gap alone do
not select the massive square-root branch.  Third, the isolated spectral
branch, its low-momentum envelope and its reconstructed quadratic action are
kept distinct from physical branch selection and from a microscopic closure
solution.  Fourth, the resulting continuous action charge identifies precisely
why $Q=1$ and the empirical scale $J_0=\hbar$ require further physical input.
The scoped negative results on free localization, physical complex-structure
selection and sector/statistics identification serve the same purpose: they
exclude shortcuts, not alternative microscopic theories generally.

Thus this is a \emph{conditional reconstruction and constraint analysis}, with
a proposed C-Motion interpretation of the assumptions, not a derivation of
ordinary quantum mechanics from established ElEmEr field equations.  It does
not yet predict a deviation from the conventional one-particle theory or
identify an observation that favors the proposed ontology.  Its testable
value at this stage lies in making the missing constitutive and selection
mechanisms explicit, so that proposed native completions can be evaluated
against them.

\section{From local-\texorpdfstring{$c$}{c} propagation to a gapped effective sector}
\label{sec:antecedents}

Primitive local-$c$ propagation by itself describes open propagation, not a massive particle.
Recurrence introduces a second scale of description.  The field can keep propagating locally at
$c$ while an organized return pattern carries a repeatable internal phase and the pattern as a
whole changes more slowly.  A recurrent structure therefore motivates looking for a
characteristic internal frequency, but recurrence alone does not determine its value or even
guarantee the particular spectral form used here.

In the effective sector studied below, that internal scale appears as a nonzero spectral
\emph{gap}.  We denote the half-gap by $\Delta>0$.  Here $\Delta$ is specifically the frequency
separation between the zero-frequency reference and the two gapped spectral families produced
by the operator.  It should not be silently identified with every recurrence frequency that a
future microscopic closure might possess.  The present theory does not yet derive what
microscopic closure dynamics creates or selects $\Delta$.

To ask what a local-$c$ gapped sector is allowed to look like, we make five explicit effective
assumptions:
\begin{enumerate}[leftmargin=*,label=\textbf{A\arabic*:},ref=A\arabic*]
\item\label{ass:qf-A1} \textbf{Local linear first-order Hermitian evolution.}  In the free sector of interest, the
resolved effective field obeys a Hermitian evolution equation first order in primitive
succession $s$ and first order in spatial derivatives.
\item\label{ass:qf-A2} \textbf{Translation-invariant free region.}  The coefficients are constant where the
free symbol, spectral branches, and translation charges are constructed.
\item\label{ass:qf-A3} \textbf{Exact isotropic characteristic speed.}  Removing the gap leaves only the
isotropic characteristics $\omega=\pm c|\mathbf k|$.
\item\label{ass:qf-A4} \textbf{Unsplit zero-momentum gap.}  At $\mathbf k=0$, the Hermitian gap
symbol has only the spectral values $+\Delta$ and $-\Delta$, with $\Delta>0$
(and possible multiplicities), rather than several inequivalent gap magnitudes.
\item\label{ass:qf-A5} \textbf{Even squared free symbol.}  For the constant-coefficient free
symbol $h(\mathbf k)=cA^jk_j+\Delta B$, the selected unsplit sector obeys
$h(\mathbf k)^2=h(-\mathbf k)^2$ for every $\mathbf k$.  This excludes
momentum-odd contributions to the \emph{squared} generator; it is an
independent restriction on the effective model class, not a consequence of local
$c$-propagation or a proof of microscopic parity symmetry.
\end{enumerate}

The first-order form, the gap, and the unsplit squared-symbol property are upstream
effective assumptions.  Translation invariance is an effective free-sector condition.  The
isotropic characteristic speed is how the inherited local-$c$ premise is represented inside
this operator class.  The physical meaning of $\Delta$ beyond the spectrum remains to be
supplied by closure dynamics; in particular, it is not identified here with a primitive
electroweak bare mass.  Conditions A4 and A5 are deliberately stronger than merely requiring
some isotropic nonzero gap: other first-order models may have split or shifted branches.

These restrictions pose a sharp mathematical question: which matrix algebra realizes a
local first-order evolution with exactly isotropic gapless characteristics, a single
zero-momentum gap magnitude, and an even squared symbol?  For finite-dimensional complex realizations, the algebraic answer is Clifford
structure; irreducibility is needed only for the subsequent minimality and
action-uniqueness statements.  The evenness condition is not derived
from the other four assumptions.

\section{Why Clifford structure appears}
\label{sec:clifford}

A single scalar amplitude can describe a wave magnitude, but it has too little internal
structure to represent several first-order propagation directions together with an isotropic
mixer without introducing direction-dependent artifacts.  The multicomponent field used here
provides an internal state space in which direction-sensitive propagation and direction-neutral
mixing can coexist.  Its components are effective internal coordinates; they are not identified
here with electric/magnetic fields, weak chirality, particle/antiparticle labels, or any other
empirical pair.

The role of the internal matrices is to encode that structure while keeping the characteristic
speed independent of spatial direction.  This is precisely the algebraic problem solved by a
Clifford representation.  We write the effective evolution as
\begin{equation}
 i\partial_s\Psi=H_s\Psi,
 \qquad
 H_s=-icA^j\partial_j+\Delta B,
 \label{eq:firstorder}
\end{equation}
where $\Psi$ is an $N$-component effective field and the constant matrices $A^j$ and $B$ are
Hermitian.  The matrices $A^j$ govern the directional first-order propagation; $B$ mixes the
internal components in the gapped sector.  The operator $H_s$ has units of inverse
succession interval (angular frequency).  When the universal action-normalization slot
$J_0$ is introduced below, the corresponding effective energy generator is
$H_E:=J_0H_s$; thus $iJ_0\partial_s\Psi=H_E\Psi$.  This is a generator convention,
not an identification of primitive succession with an independently measured clock time.

\begin{theorem}[Clifford reduction of the effective sector]
\label{thm:clifford}
Under A1--A5 in \Cref{sec:antecedents}, for a finite-dimensional complex
realization of the selected gapped sector, the coefficients satisfy
\begin{align}
 \{A^i,A^j\}&=2\delta^{ij}I,\label{eq:cliffA}\\
 B^2&=I,\qquad \{A^i,B\}=0.\label{eq:cliffB}
\end{align}
Consequently the momentum-space Hermitian symbol
\begin{equation}
 h(\mathbf k)=cA^j k_j+\Delta B
\end{equation}
satisfies
\begin{equation}
 \boxed{h(\mathbf k)^2=(\Delta^2+c^2|\mathbf k|^2)I},
 \label{eq:squaresymbol}
\end{equation}
and therefore has the two spectral values
\begin{equation}
 \boxed{\omega_\pm(\mathbf k)=\pm\sqrt{\Delta^2+c^2|\mathbf k|^2}}.
 \label{eq:dispersion}
\end{equation}
\end{theorem}

\begin{proof}
The gapless symbol is $h_0(\mathbf k)=cA^j k_j$.  Exact isotropic $\pm c$
characteristics require $h_0(\mathbf k)^2=c^2|\mathbf k|^2I$ for every $\mathbf k$.
Comparison of the quadratic coefficients gives \eqref{eq:cliffA}.  With the constant mixer,
\[
h(\mathbf k)^2=c^2|\mathbf k|^2I+c\Delta k_j\{A^j,B\}+\Delta^2B^2.
\]
The zero-momentum spectrum consists only of $\pm\Delta$; since $B$ is Hermitian,
this fixes $B^2=I$.  Assumption A5 makes the square equal at $\mathbf k$ and
$-\mathbf k$.  Subtracting the two displayed squared symbols gives
$2c\Delta k_j\{A^j,B\}=0$ for all $\mathbf k$, and hence
$\{A^j,B\}=0$ (as $\Delta>0$).  Equation~\eqref{eq:squaresymbol} and the
eigenvalues \eqref{eq:dispersion} follow.
\end{proof}

No step of this algebraic proof requires irreducibility. Reducible direct sums
also satisfy the relations; minimal representation dimension and the scalar
commutant used later require their separate irreducibility hypothesis.
The square-root dispersion follows algebraically within this restricted operator class;
it is not an independently postulated dispersion relation.  To see why the fifth
assumption matters, consider the Hermitian four-component example
$A^j=\sigma_j\otimes I_2$ and $B=I_2\otimes\sigma_3$, with Pauli matrices
$\sigma_j$.  It has the stipulated gapless $\pm c|\mathbf k|$ spectrum and
zero-momentum eigenvalues $\pm\Delta$, but $[A^j,B]=0$, giving the shifted
branches $\pm c|\mathbf k|\pm\Delta$, not \eqref{eq:dispersion}.
Thus the unsplit squared-symbol condition is a separate effective-model assumption.  Merely supplying local
first-order isotropic propagation and a nonzero gap does \emph{not} force this particular
massive-looking branch.  C-Motion has not derived the microscopic first-order operator,
the gap mechanism, or the physical selection of this unsplit sector.

\subsection{What the Clifford structure gives---and what it does not}

The smallest complex gapless block has dimension two, while the smallest block admitting an
independent anticommuting gap mixer has dimension four
\cite{LawsonMichelsohn1989SpinGeometry}.  This gives a natural spinorial container for the
gapped branch and fixes its dispersion, but it does not yet make a localized particle.  The
translation-invariant free Clifford operator has no nonzero square-integrable stationary
eigenstate: at fixed energy its Fourier support lies on a measure-zero shell.  Some additional
interaction or closure dynamics is therefore required to localize the branch physically.

The same representation theory also supplies a clean exclusion.  Irreducible complex
Clifford/Weyl module dimensions occur in powers of two, so the observed three fermion families
cannot simply be the dimension of one such module.  The result says where the family
multiplicity does \emph{not} come from; it does not supply a replacement family mechanism.
Likewise the spectral gap $\Delta$ is not, by itself, a derivation of weak chirality, a
scalar/Higgs mechanism, or a primitive gauge-invariant charged-lepton mass.

\section{From the gapped branch to the Schrödinger envelope}
\label{sec:schbridge}

The gapped spectrum contains two separated frequency families.  A \emph{positive-frequency
branch} is the family with $\omega=+\sqrt{\Delta^2+c^2k^2}$; choosing it means restricting the
effective state to that spectral subspace.  Mathematically this selects one sense of phase
evolution in the gapped sector.  The projector performs that restriction cleanly, but it is not
itself a physical law explaining why nature should occupy that branch.

Once the branch is chosen, two scales become visible.  The factor set by $\Delta$ advances
rapidly even for very small spatial momentum, whereas the extra phase associated with spatial
modulation changes slowly when $c|\mathbf k|\ll\Delta$.  This is the familiar
carrier--envelope separation used for modulated waves: the fast factor and the slow envelope
are parts of one field state.  Here that separation is what turns the square-root spectrum into
a low-energy one-particle envelope dynamics.

Define
\begin{equation}
 \Omega_k=\sqrt{\Delta^2+c^2|\mathbf k|^2},
 \qquad
 P_\pm(\mathbf k)=\frac12\left(I\pm\frac{h(\mathbf k)}{\Omega_k}\right).
\end{equation}
For data entirely in one spectral branch, $P_\pm\widehat\Psi=\widehat\Psi$, the exact
branch evolution is multiplication by $\pm\Omega_k$.  On the positive branch,
\begin{equation}
 i\partial_s\Psi_+=\sqrt{\Delta^2-c^2\nabla^2}\,\Psi_+.
 \label{eq:sqrtfull}
\end{equation}
Here and below the square root is the positive self-adjoint Fourier multiplier on the chosen
branch.  The spectral projector cleanly describes the branch once chosen; the microscopic
reason nature should select the corresponding positive-frequency/positive-energy solution
space remains unresolved.

\paragraph{Local full equation versus nonlocal isolated branch.}
The full free first-order equation \eqref{eq:firstorder} is a \emph{local} differential
system with characteristic speeds $\pm c$ under the stated effective assumptions.
By contrast, $P_\pm(\mathbf k)$ depends on $\mathbf k$ and projects the complete field
by a spatially nonlocal Fourier multiplier; the scalar multiplier
$\sqrt{\Delta^2-c^2\nabla^2}$ in \eqref{eq:sqrtfull} is likewise nonlocal as a
spatial operator.  Thus a strictly localized full-field preparation is not in
general preserved as a strictly localized \emph{positive-branch component} after
projection.  The group velocity
$|\nabla_{\mathbf k}\Omega_k|=c^2|\mathbf k|/\Omega_k\leq c$ does
not by itself establish compact-support propagation for the projected variable.
This mathematical nonlocality is a consequence of resolving one spectral sector
of a local parent description; it is \emph{not}, without further physical input,
a mechanism for superluminal signalling or a failure of causal propagation in the
complete field.  Physical preparation, interactions, and local observables cannot
be inferred from a branch projection alone.  The native \ElEmEr{} dynamics and any
physical rule selecting or sustaining this branch remain Open.

Hegerfeldt's localization obstruction concerns positive-energy quantum
evolution together with sharp particle-localization assumptions, including
orthogonality of states localized in disjoint regions
\cite{Hegerfeldt1974}. It is not a theorem that every spatially nonlocal
operator permits operational acausal signals. Its localization assumptions
must be checked for the proposed physical observables; a projected field
component alone supplies neither a localization measurement nor a signal protocol.

Finite-speed causal evolution of the complete effective field, operational
no-signalling of detector marginals, and \emph{operator microcausality} are distinct
conditions.  In particular, neither the locality of the full classical/free
operator nor a bound on the isolated-branch group velocity establishes spacelike
commutation of physical extended-mode observables.  Such an operator algebra and
its commutation properties remain Open; no spin--statistics conclusion follows
from the free branch or its spectral projection alone.

Now separate the rapid internal phase from the slowly varying spatial modulation:
\begin{equation}
 \Psi_+(\mathbf x,s)=e^{-i\Delta s}\psi(\mathbf x,s).
\end{equation}
The field $\psi$ is the \emph{envelope field}: it tracks how the amplitude and phase of the
selected branch vary on scales slow compared with the internal frequency $\Delta$.  A useful
picture is a rapidly repeating internal rhythm whose overall modulation drifts only slightly
from cycle to cycle.  The picture is an interpretation of the scale separation, not a claim
that a microscopic closure solution with that detailed motion has already been constructed.
Then
\begin{equation}
 \boxed{
 i\partial_s\psi=
 \left[\sqrt{\Delta^2-c^2\nabla^2}-\Delta\right]\psi.}
 \label{eq:exactsqrt}
\end{equation}
This is exact within the selected free branch.

For a wavepacket whose Fourier support obeys
\begin{equation}
 \epsilon=\frac{c|k|}{\Delta}\ll1,
 \label{eq:epsilon}
\end{equation}
the multiplier has the controlled expansion
\begin{align}
 \sqrt{\Delta^2+c^2k^2}-\Delta
 &=\frac{c^2k^2}{2\Delta}-\frac{c^4k^4}{8\Delta^3}
 +\Order(\epsilon^6\Delta),
\end{align}
and therefore
\begin{equation}
 \boxed{
 i\partial_s\psi=
 -\frac{c^2}{2\Delta}\nabla^2\psi
 -\frac{c^4}{8\Delta^3}\nabla^4\psi+\ldots.}
 \label{eq:slowplain}
\end{equation}
Introduce the symbolic action normalization $J_0$ and define the corresponding branch mass
scale
\begin{equation}
 \boxed{M=\frac{J_0\Delta}{c^2}}.
 \label{eq:Mdef}
\end{equation}
Multiplying \eqref{eq:exactsqrt} by $J_0$ gives the exact form
\begin{equation}
 iJ_0\partial_s\psi=
 \left[\sqrt{M^2c^4-J_0^2c^2\nabla^2}-Mc^2\right]\psi,
 \label{eq:exactsqrtJ}
\end{equation}
and the slow-envelope reduction becomes
\begin{equation}
 \boxed{
 iJ_0\partial_s\psi=
 -\frac{J_0^2}{2M}\nabla^2\psi
 -\frac{J_0^4}{8M^3c^2}\nabla^4\psi+\ldots.}
 \label{eq:schbridge}
\end{equation}
The leading term is the standard free Schr\"odinger kinetic operator with the action scale
$J_0$ left symbolic.  The result is controlled by the small parameter
$c|\mathbf k|/\Delta$: the Schrödinger equation appears here as the slow-envelope limit of the
gapped branch rather than as the primitive evolution law.  The next term also fixes the first
relativistic correction within the same expansion.

\paragraph{Physical meaning of the mass scale.}
Equation~\eqref{eq:Mdef} makes the central physical distinction explicit.  The resolved
primitive propagation tangent can remain locally at $c$ even when the recurrent organization
has a subluminal collective or envelope motion.  Mass-like behaviour belongs to the organized
pattern, not to a primitive carrier that has slowed below $c$.  Once the action scale $J_0$ is
specified, the gap carries the energy scale $J_0\Delta$ and therefore the effective mass scale
$M=J_0\Delta/c^2$.  For a branch carrying action charge $Q$, the rest-frequency contribution
is $J_0Q\Delta=QMc^2$.  The $M$ in the Schr\"odinger kinetic operator is
therefore a \emph{branch mass parameter per unit action charge}; it is not
by itself the total mass of an arbitrary finite-$Q$ field configuration.
For a single monochromatic mode with total translation momentum $\mathbf P$,
the low-momentum collective kinetic term has denominator $2QM$, not $2M$.
For a finite packet there is also energy from its nonzero Fourier-width, as
shown explicitly in \eqref{eq:packetmassdecomp} below.  This gives a
conditional effective realization of rest-like energy and collective inertia
without establishing the mass of a localized physical closure.  The physical
identification still requires native gap and localization dynamics, the
selection of $Q$, and the relevant energy--momentum/clock bridge.

\paragraph{Complex notation and physical complex structure.}
The slow envelope can be rewritten in two real quadratures with a real operator $\mathbb J$
obeying $\mathbb J^2=-I$; multiplication by $i$ is then a quarter-turn in that real
quadrature plane.  Complex notation is therefore mathematically available without having yet
explained why one particular complex structure is the physical one shared by free, bound, and
interacting states.  That selection problem is returned to in
\Cref{sec:complexmeasurement}.

\section{Phase symmetry, action amount, energy, and momentum}
\label{sec:actionbridge}

The previous section derived an evolution equation.  Its phase already carries a simple
physical meaning: it records how the effective mode advances through its recurrent oscillation,
while phase differences distinguish interference and relative organization.  To turn that
phase evolution into physical energy and momentum one needs more than periodicity; one needs an
action whose continuous symmetries generate conserved quantities.

In the present effective sector the first-order Hermitian branch admits such a local quadratic
action.  The factor $J_0$ sets its overall action scale.  Changing $J_0$ rescales the action and
all associated dimensionful Noether charges without changing the classical field equation.
This is why the branch can determine the structure of the energy--frequency relation while
leaving its absolute physical normalization unresolved.

A manifestly real representative is
\begin{equation}
 S_{\rm eff}=J_0\int\dd s\,\dd^3x\left[
 \frac{i}{2}\left(\Psi^\dagger\partial_s\Psi
 -(\partial_s\Psi^\dagger)\Psi\right)
 -\Psi^\dagger H_s\Psi\right],
 \label{eq:Seffreal}
\end{equation}
which differs by a boundary term from
\begin{equation}
 \boxed{
 S_{\rm eff}=J_0\int\dd s\,\dd^3x\,
 \Psi^\dagger(i\partial_s-H_s)\Psi.}
 \label{eq:Seff}
\end{equation}
Varying $\Psi^\dagger$ gives the same effective first-order equation.  The construction is
therefore an action for the resolved branch dynamics.  It should be read at that scale: the
microscopic \ElEmEr{} action from which such a branch might ultimately emerge is still
unknown.

\subsection{How much of the effective action is fixed}

Suppose a same-field local quadratic first-order action uses a constant Hermitian symplectic
metric $G$ and is required to reproduce the same Hermitian minimal irreducible Clifford
dynamics without changing the generator class.  Compatibility requires $G$ to commute with
all $A^j$ and $B$.  For the minimal irreducible complex module, Schur's lemma makes $G$
scalar.  Within this restricted class, the action is therefore fixed up to the common overall
normalization $J_0$ and boundary terms.  Reducible sectors can have a larger commutant and need
not share this uniqueness.  The algebra fixes the form much more tightly than it fixes the
absolute size of the action.

\subsection{Global phase symmetry and action charge}

Within the effective branch, the absolute origin chosen for the common phase is arbitrary:
shifting every component by the same constant angle does not change the action, whereas
relative phases between components or modes can affect interference and other relational
observables.  The exact global transformation
$\Psi\mapsto e^{i\alpha}\Psi$ therefore leaves the effective action unchanged.  Noether's
theorem associates a conserved quantity with this continuous phase symmetry:
\begin{equation}
 \boxed{Q=\int \Psi^\dagger\Psi\,\dd^3x.}
 \label{eq:Q}
\end{equation}
We call $Q$ the \emph{action charge}.  Here ``charge'' has the Noether meaning: it is a
conserved quantity associated with an exact continuous symmetry.  It is not electric charge,
a local gauge charge, or automatically a particle number.  In the chosen normalization $Q$
is dimensionless, and the corresponding dimensionful action amount is $J_0Q$.

Because the effective theory is classical and linear, its amplitude can be scaled
continuously, so $Q$ is continuous.  The symmetry therefore identifies where an action quantum
would have to enter, but the quadratic branch theory does not itself select a discrete unit
such as $Q=1$.

\paragraph{Three distinct quantization questions.}
It is useful to separate three logically independent questions that can otherwise be
blurred together.  First, closure dynamics, return conditions, or boundary conditions
may select a discrete set of admissible states, frequencies, or spatial eigenmodes.
This is \emph{state or mode selection}.  Second, even within one allowed mode, the
present linear effective theory permits continuous amplitude rescaling and therefore
continuous action charge $Q$; a further physical principle would be required to select
discrete action amounts or an elementary value such as $Q=1$.  Third, even if such a
unit were physically selected, its dimensional size would still be set by the
independent action-normalization scale $J_0$.  Recovering ordinary quantum
normalization would then additionally require the empirical identification
$J_0=\hbar$.

Thus
\[
\boxed{
\begin{gathered}
\text{discrete closure-state selection}
\;\neq\;
\text{action-charge quantization}
\\[2pt]
\neq\;
\text{absolute action normalization}
\end{gathered}
}
\]
and success at any one of these levels does not by itself solve the others.  In
particular, a discrete spectrum does not derive $Q=1$, and a physical selection
$Q=1$ would still not derive $J_0=\hbar$.

\subsection{Translations give energy and momentum}

For fields with the usual decay or boundary conditions, $s$-translation and spatial
translation give
\begin{align}
 E&=J_0\int\Psi^\dagger H_s\Psi\,\dd^3x,
 \label{eq:Echarge}\\
 \mathbf P&=J_0\int\Psi^\dagger(-i\nabla)\Psi\,\dd^3x.
 \label{eq:Pcharge}
\end{align}
The same action therefore ties phase evolution, spatial wavevector, energy, and momentum to
one symmetry structure.  The familiar generators $H_s$ and $-i\nabla$ are not being added here
as independent quantum postulates; they arise as the translation generators of the effective
action.

For a resolved monochromatic plane-wave branch mode
$\Psi=u\,a\,e^{i\mathbf k\cdot\mathbf x-i\omega s}$, normalized over a finite box or as a
wavepacket limit so that \eqref{eq:Q} is finite, the branch eigenvalue equation gives
\begin{equation}
 \boxed{E=J_0Q\omega,
 \qquad
 \mathbf P=J_0Q\mathbf k.}
 \label{eq:modeEP}
\end{equation}
For a single resolved mode, frequency and wavevector therefore carry energy and momentum in
proportion to the same action amount $J_0Q$.  Equivalently,
\[
 \frac{E}{\omega}=J_0Q,
 \qquad
 \frac{E}{f}=2\pi J_0Q
 \quad (f=\omega/2\pi).
\]
This makes precise, in a controlled setting, the physical intuition that converting a cycle
rate into an energy requires an action scale: $J_0Q$ is the action amount per radian of phase
advance, and $2\pi J_0Q$ is the corresponding energy--period product for one full cycle.  The
result is exact inside this linear sector and does not require choosing the special
normalization $Q=1$.

For an arbitrary finite packet in one isolated branch, let $a(\mathbf k)$ denote the spectral
amplitude in an orthonormal branch basis and define
\begin{equation}
 Q=\int |a(\mathbf k)|^2\,\dd^3k,
 \qquad
 \langle F\rangle_Q=\frac{1}{Q}\int F(\mathbf k)|a(\mathbf k)|^2\,\dd^3k.
\end{equation}
Then the exact charges are
\begin{equation}
 \boxed{
 E=J_0Q\langle\omega\rangle_Q,
 \qquad
 \mathbf P=J_0Q\langle\mathbf k\rangle_Q.}
 \label{eq:packetEP}
\end{equation}
A finite packet is therefore not assigned one exact $(\omega,\mathbf k)$ unless its spectral
support justifies that approximation.  After rest-frequency removal on the positive branch,
the corresponding envelope energy is
\begin{equation}
 E_{\rm env}=J_0Q\left\langle\sqrt{\Delta^2+c^2k^2}-\Delta\right\rangle_Q.
\end{equation}

\paragraph{Finite-$Q$ collective motion versus spectral-width energy.}
For a positive-branch packet with $Q>0$ and spectral support in
$c|\mathbf k|/\Delta\ll1$, define
$\sigma_k^2:=\langle|\mathbf k|^2\rangle_Q-|\langle\mathbf k\rangle_Q|^2$
(the sum of the component variances).  Using \eqref{eq:packetEP} and
\eqref{eq:Mdef}, the low-momentum expansion gives
\begin{equation}
\begin{split}
 E&=QMc^2+\frac{J_0^2Q}{2M}\langle|\mathbf k|^2\rangle_Q
      +\Order(J_0Q\Delta\epsilon_{\max}^4)\\
  &=QMc^2+\frac{|\mathbf P|^2}{2QM}
    +\frac{J_0^2Q}{2M}\sigma_k^2
    +\Order(J_0Q\Delta\epsilon_{\max}^4),
\end{split}
\label{eq:packetmassdecomp}
\end{equation}
where $\epsilon_{\max}$ bounds $c|\mathbf k|/\Delta$ on the packet's
spectral support.  In the finite-box monochromatic limit $\sigma_k^2=0$,
the charge-weighted energy has the usual low-momentum collective term
$|\mathbf P|^2/(2QM)$.  For a \emph{localized} packet,
$\sigma_k^2>0$: the extra energy is not part of the motion of the spectral
mean $\langle\mathbf k\rangle_Q$.  Consequently $QM$ must not be called
the full invariant mass of a generic packet or a derived physical particle
mass.  The distinction also does not turn $Q$ into a particle count.

\subsection{Why the simple product law works here}

The result above also explains when the simple energy--frequency product is legitimate.  The
Main Framework gives the Conditional stationary-family action--charge derivative relation
\begin{equation}
 \frac{\dd E}{\dd Q}=J_0\Omega
 \label{eq:conjugacy}
\end{equation}
for a smooth family of stationary points of $E-J_0\Omega Q$, with exact
phase symmetry, fixed action, couplings, boundary conditions and other
charges, vanishing boundary variations, and $Q$ a valid local family coordinate
\cite{CmotionMainFramework2026}.  In the present linear resolved-mode sector the frequency is
amplitude independent, so varying the mode amplitude changes $Q$ at fixed $\omega$ and
\eqref{eq:modeEP} indeed gives
\begin{equation}
 \frac{\dd E}{\dd Q}=J_0\omega,
 \qquad E=J_0Q\omega
\end{equation}
after the natural zero-energy reference of that linear mode.  In a nonlinear closure, $\Omega$
can vary with $Q$, so the general relation is an integral over the family rather than this simple
product.  The resolved-mode formula is
therefore a real structural result without being a universal nonlinear closure law.

\section{Where the quantum action scale enters}
\label{sec:normalization}

The effective action now tells us exactly what has and has not been fixed.  Its equations determine
how one common action scale propagates into energy and momentum, but a classical rescaling of
the whole action leaves the equations of motion unchanged.  The symbol $J_0$ denotes that
overall \emph{action-normalization scale}.  It converts the dimensionless phase/action
structure of the branch into dimensionful action, energy, and momentum.

The hierarchy is
\begin{equation}
 \boxed{
 \begin{array}{rcl}
 J_0 &:& \text{structural common action-normalization slot},\\
 Q &:& \text{continuous effective branch charge},\\
 J_0Q &:& \text{dimensionful effective action amount},\\
 Q=1 &:& \text{normalization convention unless physically selected},\\
 J_0=\hbar &:& \text{empirical identification; not fixed by the present dynamics}.
 \end{array}}
 \label{eq:hierarchy}
\end{equation}
The hierarchy says more than ``a constant is missing.''  Phase supplies a dimensionless measure
of cycle advance, while energy and momentum are dimensionful.  An action scale is exactly what
converts one into the other.  The effective branch therefore locates the junction at which the
empirical quantum of action must enter.  What it does not yet explain is whether the dynamics
selects a discrete value of $Q$, fixes the common scale $J_0$, or supplies both ingredients.
Neither happens in the present quadratic branch theory.

Consequently, for a normalized $Q=1$ resolved mode,
\begin{equation}
 E=J_0\omega,\qquad \mathbf p=J_0\mathbf k,
 \label{eq:Jrelations}
\end{equation}
and only after the empirical identification $J_0=\hbar$ does one obtain
\begin{equation}
 E=\hbar\omega=hf,
 \qquad
 \mathbf p=\hbar\mathbf k.
 \label{eq:empiricalPlanck}
\end{equation}
Identifying $J_0$ with Planck's reduced constant would connect the effective branch directly
to empirical quantum normalization.  At present that identification is supplied from
experiment rather than derived from the microscopic C-Motion dynamics.  Closure periodicity,
parameter independence of dimensionless quantities along an effective-model
scale-similarity orbit, and the Conditional recurrent conjugacy
\eqref{eq:conjugacy} do not fix it.  Likewise, normalizing one branch to $Q=1$ does not turn a
continuous classical Noether charge into a physically quantized one.

The action-current statement proved here also remains specific to the reconstructed effective
branch.  A unique positive action current for an arbitrary finite microscopic \ElEmEr{} wave
has not yet been established. The generic finite-wave action-current provenance,
its absolute normalization and its persistence through nonlinear closure
formation remain Open (sharp), independently of the resolved-branch formulas
and the empirical assignment $J_0=\hbar$ \cite{CmotionMainFramework2026}.

\section{What follows in the low-energy one-particle regime}
\label{sec:controls}

Once the gapped branch has reduced to the Schrödinger envelope and its translation generators
have been identified, familiar one-particle quantum mathematics becomes available.  These
results test whether the effective C-Motion construction has entered the right low-energy
regime.  Their physical content is easier to read in that order: first obtain the branch and
its generators, then ask whether ordinary packet localization, interference, barrier
penetration, and bound states follow.

\subsection{Wave localization and momentum uncertainty}
\label{sec:uncertainty}

For an $L^2$ packet with finite position and wavevector variances in one
Cartesian direction,
\begin{equation}
 \Delta x\,\Delta k\geq\frac12
 \label{eq:fourier}
\end{equation}
is the exact Fourier/Cauchy--Schwarz spread theorem. The established
uncertainty results provide the corresponding quantum benchmark
\cite{Heisenberg1927,Kennard1927,Robertson1929}.  Physically, making a packet narrow in
space requires combining a wider range of spatial wavelengths; a nearly single-wavelength
state necessarily extends farther.  No measurement disturbance is needed for this kinematic
trade-off.  The branch translation theorem now supplies the conditional momentum
scale.  For a fixed packet of nonzero effective charge $Q$, normalize the
spectral density as
$w(\mathbf k)=|a(\mathbf k)|^2/Q$ and take $\Delta k$ to be its standard deviation in the
chosen Cartesian component.  The conserved Noether momentum of this \emph{whole}
configuration is the definite number $P_{\mathrm{total}}=J_0Q\langle k\rangle_Q$
from \eqref{eq:packetEP}; it has no Fourier-component variance as a fixed classical
field charge.  Instead define the conditional \emph{spectral component-momentum}
variable $p_{\mathrm{sp}}(k)=J_0k$.  Its spread satisfies
\begin{equation}
 \boxed{\Delta x\,\Delta p_{\mathrm{sp}}\geq\frac{J_0}{2}.}
 \label{eq:uncertaintysp}
\end{equation}
One may alternatively introduce the explicitly charge-scaled spectral variable
$P^{(Q)}_{\mathrm{sp}}(k)=J_0Qk$, which gives the bookkeeping identity
\begin{equation}
 \boxed{\Delta x\,\Delta P^{(Q)}_{\mathrm{sp}}\geq\frac{J_0Q}{2}.}
 \label{eq:uncertaintyQ}
\end{equation}
The second spread is not an uncertainty in the integrated field momentum
$P_{\mathrm{total}}$.  For an explicitly normalized $Q=1$ branch the two
spectral variables coincide.  Only with the empirical identification $J_0=\hbar$
does their coefficient become $\hbar/2$.  What has been obtained is a
Fourier/translation-component relation, not yet the momentum measurement-outcome
uncertainty of a physical single particle; that interpretation requires a further
quantum measurement framework.

\subsection{Superposition}

Linearity of the Clifford equation, the isolated branch generator, and the slow envelope means
that linear combinations of solutions are again solutions while the system remains in this
free linear sector.  In a field-first picture this is less mysterious than the familiar phrase
\emph{a particle in two places}: the effective state is one spatially distributed field
configuration, and different path or mode contributions can coexist and interfere within it
\cite{Arndt1999C60,Fein2019superposition}. Relative phase determines the cross-terms.
These experiments demonstrate interference of composite systems; they do not
uniquely establish the C-Motion field ontology or its microscopic dynamics.

This statement belongs to the effective propagation regime.  It does not require the complete
microscopic \ElEmEr{} dynamics to be globally linear, and a later interaction may reorganize
the complete local field state outside this reduction.  Linearity also does not, by itself,
assign Born probabilities to the amplitudes.

\subsection{Tunnelling}

The dependency chain is
\begin{equation}
\begin{gathered}
 \text{Clifford sector under the assumptions above}
 \longrightarrow
 \text{slow-envelope kinetic operator}\\
 \longrightarrow
 \text{standard barrier mathematics once }V\text{ is separately supplied}.
\end{gathered}
\end{equation}
Once a suitable \emph{low-energy interacting envelope} and a potential $V$ are separately
supplied, the usual Schr\"odinger barrier equation supports evanescent penetration through
classically forbidden regions.  This is not an exact theorem that an arbitrary local $V$
can be added to the \emph{free} positive-frequency projector: an inhomogeneous interaction
can mix the two free spectral branches.  The field need not be pictured as a point particle
traversing a region with locally forbidden classical kinetic energy.  Instead, the effective
envelope is exponentially small within a barrier and can remain nonzero beyond it.
The standard WKB exponent is recorded in \Cref{app:compatibility}
\cite{Wentzel1926WKB}.  The C-Motion-specific part of this conditional chain is the free
kinetic-envelope antecedent; the interaction, its admissible projection/reduction, the
source, sign and coupling of $V$, and the physical transmission-rate interpretation
remain additional assumptions or open native problems.

\subsection{Hydrogen}

Hydrogen is included as a downstream \emph{conditional consistency test}, not as a new
derivation of atomic physics from C-Motion primitives.  The dependency chain is:
\begin{enumerate}[leftmargin=*]
\item the Clifford free sector under all five effective assumptions of
  \Cref{sec:antecedents}, including the independently stipulated even-squared-symbol condition;
\item its controlled free slow-envelope kinetic operator \eqref{eq:schbridge} and the
  symbolic action normalization $J_0$;
\item an independently supplied charged-closure/Maxwell--Coulomb interaction sector,
  including the source/current and the coupling convention
  \cite{CmotionSMStructureI2026};
\item an additional controlled low-energy \emph{interacting} one-particle/two-body reduction,
  including a supplied $Q=1$ single-particle identification (not a derived physical
  action quantum), and the standard electron--proton reduced mass $\mu$; and
\item ordinary Coulomb spectral and boundary-condition mathematics.
\end{enumerate}

\paragraph{Interaction and projection qualification.}
The exact square-root operator \eqref{eq:exactsqrt} and the free projector
$P_+^0:=P_+$ defined above are constructed in a \emph{free},
translation-invariant effective sector.  For a
spatially varying local potential $V(\mathbf x)$ one generally has
\begin{equation}
 [P_+^0,V(\mathbf x)]\ne 0,
 \label{eq:freeprojpotential}
\end{equation}
so the potential may mix the free positive- and negative-frequency subspaces.
The expression $P_+^0 V P_+^0$ is itself a projected, generally nonlocal interaction,
not identically a local multiplicative potential acting in a free one-branch Hilbert
space.  An interacting low-energy Schr\"odinger description is consequently a
\emph{further effective reduction}, rather than an exact algebraic consequence of the
free spectral projector.  In the weakly bound, low-momentum regime, with
$|E_{\rm bind}|\ll\mu c^2$, typical relative momenta
$|\mathbf p|\ll\mu c$, and small effective Coulomb coupling
$\alpha_{\rm eff}\ll1$ as defined below, standard block reduction permits the leading Coulomb--Schr\"odinger form; omitted
higher-order corrections include relativistic and spin-dependent terms.  These
conditions do not derive the electromagnetic source, interaction or the physical
branch selector from ElEmEr dynamics.

\paragraph{Coulomb coupling and its units.}
Write $g_C>0$ for the coefficient of the attractive potential $-g_C/r$; it has
units of energy times length.  For ordinary opposite charges of magnitude $e$ in
vacuum, the externally normalized SI value is
$g_C=e^2/(4\pi\varepsilon_0)$.  In this effective branch one may define the
dimensionless ratio
\begin{equation}
 \alpha_{\rm eff}:=\frac{g_C}{J_0c}.
 \label{eq:alphaeff}
\end{equation}
Only with the empirical identification $J_0=\hbar$ and the standard Maxwell
charge normalization does $\alpha_{\rm eff}$ equal the measured
fine-structure constant $\alpha_{\rm fs}$.  None of these physical identifications
follows from the free Clifford symbol alone.

Under the stated two-body, interaction-reduction and normalization assumptions,
the ordinary leading Coulomb envelope equation is
\begin{equation}
 \left[-\frac{J_0^2}{2\mu}\nabla^2-\frac{g_C}{r}\right]\psi=E\psi,
 \label{eq:hydrogenJ}
\end{equation}
and regularity and square-integrability give its gross levels,
\begin{equation}
 E_n=-\frac{\mu g_C^2}{2J_0^2}\frac{1}{n^2}
     =-\frac{\mu c^2\alpha_{\rm eff}^2}{2n^2},
 \qquad n=1,2,\ldots.
 \label{eq:hydrogenlevelsJ}
\end{equation}
The $E_n$ are binding energies relative to the leading two-body rest-energy
reference; this is not an absolute ElEmEr closure-energy prediction.
With $J_0=\hbar$, imported observed charges and the physical reduced mass,
the leading coefficient agrees with the standard nonrelativistic gross
hydrogen spectrum.  The allowed integer $n$ comes from ordinary wave
regularity, normalizability and boundary conditions even when the amplitude
remains continuous; it neither derives physical $Q=1$ nor fixes $J_0$.

This is a conditional \emph{compatibility} of the free kinetic
bridge with an additional interacting low-energy construction.  It does not
supply native charged-closure currents, the native unbroken electromagnetic $\mathrm U(1)$ current and
Maxwell source/action realization, or an exact
interacting positive-frequency selector. The separate, more ambitious problem
of deriving the surviving electromagnetic direction from electroweak structure
is not needed for this leading hydrogen control.  Static Coulomb is a
constraint/source sector in the conditional electromagnetic description, not
an on-shell open photon mode.  The radial mathematics is
summarized in \Cref{app:compatibility}.

\section{Spinorial return: what rotation structure is already available}
\label{sec:spin}

Ordinary orbital structure is not enough to represent half-integer spin.  A rotating scalar
pattern or ordinary vector returns to the same representative after a $2\pi$ rotation.  A
spinor instead transforms on the double cover of the spatial rotation group: one $2\pi$
rotation changes the representative by a sign, and only $4\pi$ returns it completely.

This is physically interesting in C-Motion because a recurrent organization may carry more
than spatial support; it may also carry an internal frame or orientation state that must return
consistently when the whole organization is rotated.  The spinorial picture is therefore not a
tiny object literally spinning in space.  It is a statement about how the internal state
transforms under rotation.  The Clifford sector supplies an appropriate mathematical container
for that possibility.

If a physical internal closure state is selected to transform in the double cover of
$\SO$, then
\begin{equation}
 U(\hat{\mathbf n},\theta)
 =\exp\!\left[-\frac{i\theta}{2}\hat{\mathbf n}\cdot\bm\sigma\right],
 \qquad U(2\pi)=-I,\quad U(4\pi)=I.
 \label{eq:4pi}
\end{equation}
The equations therefore admit the familiar spinor return law
\cite{RauchWerner2000}. Neutron-interferometry experiments tested the
relative sign under spin rotation \cite{Rauch1975Spinor,Werner1975}; they
provide the empirical spinorial benchmark, not evidence selecting a native
C-Motion closure mechanism. A natural C-Motion interpretation is that a framed recurrent
organization could carry this double-cover return structure.  What is not yet known is why a
physical elemental matter closure should be selected into that sector rather than into an
ordinary vector-like one.

Within ordinary single-valued classical scalar/vector/tensor transport, a
$2\pi$ circuit does not by itself produce the spinorial sign.  A physical
framed or protected-core configuration-space lift could carry extra return
structure, but its existence, occupation, and sector selection would be
additional native obligations.  This is a restriction on the tested ordinary
transport route, not a general exclusion of all possible classical field
organizations.

A one-body $2\pi\mapsto-1$ return does not determine the many-body exchange character.  In
three-dimensional identical-particle configuration space both $+1$ and $-1$ one-dimensional
exchange characters are topologically admissible \cite{LeinaaMyrheim1977}; a specific
Finkelstein--Rubinstein identification requires the actual physical configuration space
\cite{FinkelsteinRubinstein1968}.  Therefore the physical exchange rule
\begin{equation}
 H_{\rm exch}=-I
\end{equation}
has not been obtained from the one-body return alone.  Once antisymmetric many-body
composition is supplied, Pauli exclusion follows exactly by standard algebra.  The unresolved
physics is the bridge from a framed one-body closure to the exchange character of two
identical closures.  The present construction does not supply a C-Motion spin--statistics theorem.
The established relativistic spin--statistics argument uses additional
physical hypotheses: Pauli's free-field analysis invokes relativistic
invariance and positivity of energy, and, for its integer-spin conclusion,
commutativity of observables at spacelike separation
\cite{Pauli1940SpinStatistics}. More generally, locality/microcausality and
a positive physical state space cannot be replaced by a one-body rotation
matrix. Those field-theoretic antecedents have not been derived here.

\section{Decoherence and Zeno as consistency checks}
\label{sec:decozeno}

The effective branch can also be placed inside familiar open-system and repeated-interrogation
mathematics.  In physical terms, decoherence begins when information about relative phase is
carried into degrees of freedom that are no longer controlled as part of the system.  It is an
interaction effect, not a consequence of an observer merely ``looking.''  To calculate it one
must nevertheless specify a system--environment split and a coupling.  If such a unitary
coupling produces
\begin{equation}
 |\Psi\rangle=a|1\rangle|E_1\rangle+b|2\rangle|E_2\rangle,
\end{equation}
then tracing over the environment suppresses the system coherence by the overlap
$\langle E_2|E_1\rangle$ \cite{Zurek2003Decoherence}.  This shows compatibility with ordinary
decoherence once a microscopic environment and coupling are specified; it does not yet tell
us what the microscopic C-Motion environment, pointer basis, or coupling should be.

Similarly, a self-adjoint effective generator and a normalized initial state with finite
spectral variance give the standard quadratic short-interval survival law.  Under an
\emph{additional} ideal rank-one survival-test protocol, repeated tests with
state re-preparation conditioned on the survival result produce the standard Zeno
suppression \cite{MisraSudarshan1977}.  Such a test sequence is not obtained merely
by allowing the free \ElEmEr{} field to propagate; it supplies a measurement or
intervention rule.  The compact formula is given in \Cref{app:compatibility}.
A C-Motion-specific implementation, and especially an anti-Zeno prediction, would
require the actual apparatus/environment coupling, spectrum, correlation law, and
readout dynamics.  Neither decoherence nor ideal Zeno suppression selects a
definite single outcome or supplies the Born rule.

\section{The remaining specifically quantum ingredients}
\label{sec:complexmeasurement}

\subsection{Why the complex structure is still a physical question}

The symbol $i$ in the effective equations is easy to use and harder to explain physically.  A
real oscillation can be described by two quadratures---for example cosine-like and sine-like
components separated by a quarter cycle.  Writing those two real components as one complex
amplitude is equivalent to choosing an operator $\mathbb J$ that rotates one quadrature into
the other and satisfies $\mathbb J^2=-I$.  In that limited sense, complex notation is a compact
way to encode a two-dimensional real phase plane.

The deeper question is why one particular quarter-turn structure should be the physically
selected complex structure shared by free, bound, and interacting systems, and how that choice
is tied to a positive-frequency/positive-energy solution space.  The mathematical packaging is
available; the physical selection from the underlying real field dynamics is not yet explained.

A tested proposal identified the complex unit with the local turning map
$J_h v=v\times\widehat h$ on a nonzero-steering transverse plane.  This map
has $J_h^2=-I$ on that plane, but it does \emph{not} select the common complex
unit of the field sector.  A straight monochromatic wave accumulates eikonal
phase, with $\dd\chi/\dd\ell=k_{\rm eff}$, while its propagation direction
need not turn at all; local turning instead depends on a transverse gradient.
Literal identification also confuses independently selectable carrier-phase
winding with path/frame-turning winding, artificially collapsing the permitted
mode tower.  Thus the literal local-steering-origin identification is
\statusRetracted{} within the tested scope.  This does not exclude a future, separately derived
phase--frame coupling, and a chosen local quarter-turn or intertwiner is not a
physical selection law for the global complex structure.

A complete answer should come from the microscopic action, constraints, symplectic structure,
and positive physical norm.  Those structures are not yet available at the required level.
The branch calculations above therefore use a compatible complex representation and branch
choice without claiming to have explained why that choice is physically fundamental.
Real quadratures still permit interference.  Expressing a general closure as a
single amplitude times $e^{i\chi}$ additionally requires an independently
selected periodic, narrow-band, or otherwise phase-reducible mode description;
this is not an objection to the explicitly controlled branch/envelope calculation
above, nor a claim that arbitrary linear superpositions are monochromatic.
Native compatibility of a shared complex structure across interacting sectors,
and the physical composition rule for multiple prepared systems, are likewise
unresolved.  Neither naive real nor standard complex tensor-product composition
is presently derived from the underlying field; their experimental comparison
must not be presented as a direct prediction of this unfinished framework.

There is nevertheless a precise external composition benchmark.
Renou et al. separate real and complex Hilbert-space predictions in networks
with independently prepared sources and specified local measurements,
using ordinary tensor products for system composition
\cite{Renou2021RealQuantum}. The superconducting implementation of Chen
et al. and the optical network of Li et al. test the corresponding bounds
\cite{Chen2022RealQT,Li2022RealQT}; their conclusions retain the respective
preparation, measurement and independence assumptions. These results constrain
the standard real-Hilbert composition class, not every ontology formulated
with real fields or every alternative composition law. A real-quadrature
rewrite of a single complex system does not settle how independent systems
compose. C-Motion has not supplied that physical composition rule and therefore
does not yet confront this benchmark as a completed microscopic theory.

\subsection{Conjugate branches and the antimatter interpretation}
\label{sec:conjugate-antimatter}

The complex representation makes a conjugation operation mathematically available, but that
fact is weaker than a physical antimatter theory.  Once a complex structure has been chosen one
can write, schematically, $\Psi\mapsto\Psi^*$.  A physical charge-conjugation operation must do
more: it must map one admissible recurrent organization to another admissible physical branch
with the correct conserved quantities, interactions, and gauge-sector transformation laws.
Because the physical complex/positive-frequency structure is itself still unresolved in the
preceding subsection, ordinary algebraic complex conjugation cannot by itself be identified with
the physical $C$ operation.

A bounded C-Motion hypothesis is nevertheless available.  Particle and antiparticle states may
correspond to a pair of recurrent closure organizations related by a physical conjugate-state
map, which we may denote schematically by
\begin{equation}
 \mathcal C_{\rm cl}:\Psi_{a}\longmapsto\Psi_{\bar a}.
 \label{eq:closure-conjugate-map}
\end{equation}
The notation in \eqref{eq:closure-conjugate-map} is a candidate physical relation, not a derived
microscopic symmetry operator.  If the microscopic closure dynamics admits such a paired branch
and the two branches have the same invariant return energy, then the existing conditional
mass-as-return-energy interpretation gives equal particle and antiparticle masses.  The equality
is therefore conditional on the physical branch pairing and equal invariant return energy; it
does not follow merely from writing $\Psi^*$.

The charge statement is an additional bridge.  An opposite electric or gauge current would
follow only if the sector-specific physical current/source map changes orientation under the
conjugate-state operation.  That map is not derived in this paper.  Local gauge architecture and
electromagnetic completion belong to the Standard-Model sector \cite{CmotionSMStructureI2026}.
Accordingly, the structural conjugate-closure hypothesis should not be read as a derivation of
opposite electric charge, a completed physical $C$ symmetry, or CPT.

There is also a useful but still conjectural process interpretation.  If an interacting closure
and its physical conjugate can jointly lose the organization that keeps each branch closed, the
pair may reorganize into whatever open or other outgoing modes are allowed by the microscopic
action and conserved quantities:
\begin{equation}
 \text{closure}+\text{conjugate closure}
 \longrightarrow
 \text{allowed outgoing modes}.
 \label{eq:conjugate-unwinding}
\end{equation}
This is a candidate reading of particle--antiparticle annihilation, not a calculation of an
annihilation amplitude, allowed channel, cross section, lifetime, or branching ratio.  Those
require the still-missing interaction action and couplings.  The existence or magnitude of a
matter--antimatter asymmetry is likewise not explained by this structural pairing and remains an
independent open problem.

\subsection{Why recovering wave dynamics is not yet a measurement theory}

Recovering a Schrödinger envelope explains an important part of quantum evolution, but a
measurement must still be a physical interaction between the prepared field organization,
apparatus, and environment.  The present paper does not yet specify what microscopic
interaction makes one laboratory coupling a measurement or why individual outcomes occur with
Born probabilities.  No appeal to observation alone fills that gap.

The continuous Noether charge $Q=\int\Psi^\dagger\Psi$ does not solve the problem: choosing
$Q=1$ is a normalization of the effective branch, not a derivation of outcome weights.
Decoherence can suppress interference between effective alternatives without selecting one
actual outcome, and the Zeno effect presupposes an interrogation model rather than deriving it.

The remaining measurement problem can therefore be stated plainly: C-Motion still needs a
microscopic account of apparatus coupling, definite single outcomes, and the empirical Born
weights. The Born probability prescription is the empirical benchmark,
historically introduced for scattering \cite{Born1926}; its modern
state/effect form below is a control, not a derivation from that historical
paper or from the present effective action. Concretely, an eventual detector model must specify physically admitted
apparatus/environment preparations and their measure $\nu$, interaction and
record-selection dynamics $\mathcal S$, and show why repeated records in a
specified context have frequencies matching the standard control
$\operatorname{Tr}(\rho E_i)$ for the corresponding effective state $\rho$
and measurement effects $E_i$.  Supplying those frequencies as initial data, or
merely preserving a previously assumed measure, is not their derivation.
Historical basin-capture or threshold proposals are research candidates, not
completed physical selection laws.  The detailed task belongs to the separate,
publication-deferred measurement-completion research programme; it is not reopened here.

\section{What remains to make the bridge microscopic and fully quantum}
\label{sec:open}

The constructive chain reached in this paper can be summarized without the internal project
classification language:
\begin{equation}
\boxed{
\begin{gathered}
\text{local-$c$ field picture}\\
+\text{ stated first-order isotropic assumptions}\\
+\text{ unsplit gap/even squared symbol}\\
\Longrightarrow\text{Clifford algebra}
\Longrightarrow\omega_\pm^2=\Delta^2+c^2k^2\\
\Longrightarrow\text{selected square-root branch}
\Longrightarrow\text{slow Schrödinger envelope}\\
\Longrightarrow\text{effective branch action}
\Longrightarrow Q,E,\mathbf P\\
\Longrightarrow E=J_0Q\omega,\quad \mathbf P=J_0Q\mathbf k\\
\text{(for a resolved linear mode).}
\end{gathered}}
\label{eq:constructivespine}
\end{equation}

Several missing steps are now sharply localized.  Upstream, the microscopic C-Motion dynamics
must explain why a local first-order Clifford-like sector exists at all, how a nonzero
isotropic gap and the unsplit squared-symbol property are generated or selected, and what
microscopic action underlies the effective branch.  It must also select the physical complex/positive-frequency structure rather than
merely permit one.  The associated free-sector vacuum selection, admissible
mode/phase reduction for recurrent closures, and composition of interacting
complex state sectors remain independent native debts; interference in a real
quadrature representation is not in question.

At the normalization level, the theory must explain whether physical action comes in discrete
units, whether a physical elementary branch has $Q=1$ or some other selected value, and why
the universal action scale is $J_0=\hbar$.  Downstream, a complete particle account must supply
the relevant interaction potentials/couplings, explain why matter occupies the spinorial
sector, determine whether physical conjugate closure branches exist and what operation realizes
charge conjugation and current reversal, derive the exchange character of identical closures,
and ultimately account for Born probabilities and definite outcomes.

These questions identify the precise places where additional physics is still required.  The
present paper stops at that boundary rather than introducing new mechanisms to fill it.  The
formal scientific status of each item is collected in \Cref{app:status}.

\section{Conclusion}

The constructive result can be stated compactly.  If local-$c$ propagation admits a local
linear first-order isotropic sector with the stipulated two-valued zero-momentum gap
\emph{and} independent even-squared-symbol condition, its complex gapped operator satisfies Clifford relations. The minimal
irreducible realization is four-dimensional; this extra restriction is not
needed for the algebraic reduction.  The resulting spectrum has a rapid internal scale $\Delta$ and an exact square-root
branch.  Factoring out that rapid phase exposes a slow envelope, and expanding the branch for
$c|\mathbf k|/\Delta\ll1$ yields the Schrödinger kinetic equation with the effective mass scale
\[
 M=\frac{J_0\Delta}{c^2}.
\]
This $M$ is the resolved-branch kinetic parameter; a single mode with action
charge $Q$ carries a $Q$-weighted rest-frequency energy, while a finite packet
also retains spectral-width energy.  Neither $M$ nor $QM$ is automatically the
mass of a physical localized closure.  The physical picture is that rest-like
energy can belong to recurrent internal organization while primitive propagation
continues locally at $c$.  The positive-branch square-root evolution is a
nonlocal spectral reduction of the local effective first-order equation, not an
independent native propagation or signalling law.  The fast internal phase and
the slow envelope are two scales of one effective field organization.

The branch action then connects phase symmetry to a continuous Noether action charge and
connects translations to energy and momentum.  A resolved mode satisfies
\[
 E=J_0Q\omega,
 \qquad
 \mathbf P=J_0Q\mathbf k,
\]
while a finite packet carries the corresponding action-weighted spectral means.  The same
construction reveals exactly where empirical quantum normalization has not yet been explained:
$Q$ remains continuous in the quadratic theory, $Q=1$ is not physically selected, and the
overall action scale $J_0$ is not yet derived to equal $\hbar$.

With the low-energy free envelope in place, Fourier uncertainty and
superposition follow within that effective sector.  Tunnelling and hydrogenic
bound-state mathematics follow only with an additional low-energy interacting
reduction and supplied sector-specific potentials, not from the free positive-
frequency projection alone.  The spinorial structure
also supports the characteristic $4\pi$ return, giving a plausible container for framed
matter organization.  A further bounded structural possibility is that antiparticles are
paired physical conjugate closure branches; equal mass then follows only conditionally from
equal invariant return energy, while physical charge conjugation, opposite gauge current and
annihilation dynamics remain unresolved.  What remains genuinely quantum is now more sharply
defined rather than hidden: the physical complex/positive-frequency structure, action
quantization, exchange statistics, and the measurement dynamics that turn amplitudes into
Born-distributed definite outcomes.

The mathematical recovery of familiar free-particle formulas is a
conditional consistency result, not an independent empirical validation of the
C-Motion interpretation or a claim that the Dirac and Schr\"odinger limits
are new mathematics.  The distinct contribution is the explicitly scoped
chain of antecedents, consequences, exclusions and unresolved native
selection problems identified in \Cref{sec:contribution}.

The resulting picture is constructive but incomplete.  C-Motion supplies a concrete effective
path from recurrent local-$c$ field organization, under explicit operator and gap assumptions,
to a massive square-root branch and its Schrödinger limit.  The remaining problems are no
longer hidden inside the familiar quantum formulas: they concern the microscopic operator and
gap, the physical complex/branch selection, the quantum of action, exchange statistics, and
the interaction dynamics that produce Born-distributed definite outcomes.


\medskip

\appendix
\section{Downstream compatibility formulas}
\label{app:compatibility}

\subsection{Barrier penetration / WKB}

With a separately supplied, admissible \emph{low-energy interacting} envelope
and a one-dimensional barrier potential (not an exact generic addition to the
free spectral projector),
\begin{equation}
 -\frac{J_0^2}{2M}\psi''+V(x)\psi=E\psi,
\end{equation}
a classically forbidden region $V>E$ has local decay rate
\begin{equation}
 \kappa(x)=\frac{\sqrt{2M[V(x)-E]}}{J_0}.
\end{equation}
For a slowly varying barrier,
\begin{equation}
 T_{\rm WKB}\sim
 \exp\!\left[-2\int_{x_1}^{x_2}
 \frac{\sqrt{2M[V(x)-E]}}{J_0}\,\dd x\right].
\end{equation}
Setting $J_0=\hbar$ recovers the conventional expression
\cite{Wentzel1926WKB}.  The potential is additional sector-specific structure.

\subsection{Coulomb bound-state termination}

With the \emph{additional} weakly interacting one-particle/two-body reduction
and dimensionful coupling $g_C$, separating variables in \eqref{eq:hydrogenJ}
gives the standard Coulomb radial equation.  The regular solution becomes a
confluent hypergeometric function whose square-integrable decay requires
polynomial termination.  This gives the principal integer $n$ and
\eqref{eq:hydrogenlevelsJ}.  The calculation is standard bound-state mathematics;
it neither supplies the physical source/charge-normalization law nor proves
that an arbitrary free positive-frequency state remains in that sector after
an interaction is switched on.

\subsection{Decoherence and Zeno formulas}

For an effective system--environment state
$|\Psi\rangle=a|1\rangle|E_1\rangle+b|2\rangle|E_2\rangle$, the reduced coherence is multiplied
by $\langle E_2|E_1\rangle$ \cite{Zurek2003Decoherence}.  This assumes the factorization and
coupling.

For the effective frequency generator $H_s$ of \eqref{eq:firstorder}, define
$H_E=J_0H_s$.  Let $\delta s$ denote a finite evolution interval in the same parameter
as $H_s$ (not automatically an observer-clock duration).  For a normalized state with
finite spectral variance $\Delta_\psi H_E<\infty$, with $\delta s$ having
the units of $s$ so that $H_s\delta s$ is dimensionless,
\begin{equation}
 |\langle\psi|e^{-iH_s\,\delta s}|\psi\rangle|^2
 =|\langle\psi|e^{-iH_E\,\delta s/J_0}|\psi\rangle|^2
 =1-\frac{(\Delta_\psi H_E)^2}{J_0^2}(\delta s)^2+o((\delta s)^2).
\end{equation}
For $N$ ideal rank-one survival projections onto the same initial state,
equally spaced over a fixed interval $S$ of the \emph{same} evolution parameter,
conditional survival probabilities multiply:
\[
P_N(S)=\left[P_{\rm surv}(S/N)\right]^N\longrightarrow1
\qquad(N\longrightarrow\infty),
\]
under the finite-variance and ideal-state-reset assumptions stated above.
This is a mathematical interrogation protocol, not a native C-Motion
apparatus or outcome law \cite{MisraSudarshan1977}.  Neither the required
state-reset dynamics nor a microscopic anti-Zeno rate follows from the effective
branch alone.

\section{Separating phase symmetry, gauge redundancy, and Maxwell dynamics}
\label{app:gaugeguard}

The branch-action theorem makes an important distinction precise.  A global phase symmetry of
the specified effective action gives the Noether charge $Q$.  Declaring an overall ray-state
phase unobservable is a different statement, and neither one by itself creates a local gauge
redundancy.  If local
phase frames are introduced, a connection and transformation law are additional structures;
Maxwell dynamics additionally requires a kinetic/source law and its normalization.  Thus
\begin{align}
 \text{ray phase}&\not\Rightarrow\text{global Noether charge},\\
 \text{global action }\UU&\not\Rightarrow\text{local gauge dynamics},\\
 \text{local gauge redundancy}&\not\Rightarrow\text{Maxwell dynamics}.
\end{align}
The local electromagnetic completion is treated separately in
\emph{Gauge Architecture and Constraints in C-Motion I: Recurrent Closure,
Native Constraints, and Conditional Gauge Structure}
\cite{CmotionSMStructureI2026}.

\section{Major-result scientific-status and nativity register}
\label{app:status}

This appendix records the scientific status of the major results used in this paper.
Scientific status is separate from the route/provenance shorthand: \emph{N2} means a
calculation in the explicitly assumed C-Motion-motivated effective model class;
\emph{C} denotes a downstream standard or externally supplied mathematical/physical
control; and \emph{O} marks an open native physical input or completion.  These
codes classify how an argument enters the manuscript, not its scientific status
and not a claim that the microscopic C-Motion dynamics has been established.
Empirical calibrations and unresolved dependencies are stated in the scope column.

The scientific-status vocabulary follows the Derivation Ledger.
\emph{Conditional} denotes a conclusion from an explicit premise/model class
not itself derived; \emph{Structural} denotes an exact or banked
organizational/accounting statement that supplies no microscopic dynamical law.
\emph{Model-derived} denotes a rigorous result inside a specified effective
or analogue model, never a native Derived claim.
The two labels are not mutually exclusive as logical descriptions:
a Model-derived result may also depend on assumptions.  In this
register, Conditional emphasizes the unresolved physical premises
of a proposed C-Motion reduction, whereas Model-derived emphasizes
an exact calculation within a separately specified effective or
analogue control model.  Neither label establishes a native
microscopic derivation. \emph{Assumed} denotes a provisional active premise.
 \emph{Negative} excludes an identification, mechanism or branch
 only in its stated scope; \emph{Conjectural} is motivated
but underived. \emph{Open} is unresolved, and \emph{Open (sharp)} isolates
a missing bridge, not greater confidence. \emph{Deferred} places work outside
the publication critical path without resolving it. \emph{Retracted} records
the withdrawal of the specifically identified tested claim; it does not
exclude distinct mechanisms outside that claim's scope.

\begingroup
\small
\setlength{\tabcolsep}{3pt}
\begin{longtable}{P{3.7cm}P{2.7cm}P{1.45cm}P{6.6cm}}
\toprule
\textbf{Major result} & \textbf{Scientific status} & \textbf{Nativity} & \textbf{Scope / unresolved bridge}\\
\midrule
\endfirsthead
\toprule
\textbf{Major result} & \textbf{Scientific status} & \textbf{Nativity} & \textbf{Scope / unresolved bridge}\\
\midrule
\endhead
Local linear first-order free operator & \statusAssumed{} / \statusConditional{} & O & Required effective antecedent; microscopic origin unresolved.\\
Nonzero isotropic gap $\Delta$ & \statusAssumed{} / \statusOpen{} & O & Required for the gapped branch; physical origin/meaning unresolved.\\
A5 follows from local $c$ and a gap & \statusNegative{} & N2 & The displayed commuting $A^j,B$ counterexample satisfies A1--A4 but gives shifted branches; it excludes this implication, not all massive models or independent routes to A5.\\
Clifford anticommutation relations & \statusConditional{} & N2 & Forced within assumptions A1--A5, including the independent even-squared-symbol restriction; not from local-$c$ and a gap alone.\\
$\omega_\pm=\pm\sqrt{\Delta^2+c^2k^2}$ & \statusConditional{} & N2 & Consequence of the Clifford symbol; not inserted independently.\\
Exact isolated-branch square-root generator & \statusConditional{} & N2 & Nonlocal Fourier multiplier on the selected spectral subspace; full free first-order parent remains local. Physical branch selector and observable causal bridge Open.\\
Controlled Schr\"odinger kinetic envelope & \statusConditional{} & N2 & $c|k|/\Delta\ll1$ expansion; $J_0$ remains symbolic.\\
$M=J_0\Delta/c^2$ branch mass scale & \statusStructural{} / \statusConditional{} & N2 & Parameter per unit $Q$; monochromatic rest-frequency energy $QMc^2$. Finite packets also have spectral-width energy; localized physical mass remains unselected.\\
Quadratic effective branch action & \statusConditional{} & N2 & Reconstructed effective action, not microscopic \ElEmEr{} action.\\
Minimal irreducible action uniqueness & \statusConditional{} & N2 & Same-field local quadratic class; scalar commutant only; reducible sectors excluded from uniqueness claim.\\
$Q=\int\Psi^\dagger\Psi$ & \statusConditional{} & N2 & Continuous classical effective Noether charge.\\
$E$ and $\mathbf P$ translation charges & \statusConditional{} & N2 & Noether charges of the reconstructed effective action.\\
$E=J_0Q\omega$, $\mathbf P=J_0Q\mathbf k$ & \statusConditional{} & N2 & Exact for a resolved linear monochromatic mode; not generic nonlinear closure law.\\
Packet spectral-mean $E,\mathbf P$ & \statusConditional{} & N2 & Exact in one isolated branch; the controlled low-momentum $E$ separates collective $|\mathbf P|^2/(2QM)$ and nonzero spectral-width energy.\\
Generic nonlinear $E=J_0\Omega Q$ & \statusNegative{} & N2 & Only special linear/isochronous product; the Conditional Main Framework stationary-family relation is $\dd E/\dd Q=J_0\Omega$ under its fixed-action and exact-symmetry hypotheses.\\
$J_0$ normalization slot & \statusStructural{} & N2 & Overall action normalization remains free in the effective action.\\
Physical $Q=1$ / elementary action increment & \statusOpen{} & O & Unit normalization is not physical quantization.\\
$J_0=\hbar$ & \statusAssumed{} / \statusSharpOpen{} & O & Empirical identification; no native derivation retained.\\
$E=\hbar\omega=hf$ & \statusSharpOpen{} & C & Used downstream only through the empirical $J_0=\hbar$ normalization; a native universal energy--frequency/action-quantum derivation remains Open (sharp).\\
Fourier $\Delta x\Delta k\geq1/2$ & \statusStructural{} & C & Exact standard Fourier mathematics for an $L^2$ packet with finite variances; no physical momentum or measurement identification is included in this row.\\
Spectral spreads $\Delta x\Delta p_{\rm sp}\geq J_0/2$ and $\Delta x\Delta P^{(Q)}_{\rm sp}\geq J_0Q/2$ & \statusConditional{} & N2 & Fourier theorem plus N2 translation component-momentum identification; neither $\Delta p_{\rm sp}$ nor $\Delta P^{(Q)}_{\rm sp}$ is a variance of the fixed total-field Noether momentum.\\
Linear superposition & \statusConditional{} & C & Downstream consequence of the N2 linear sector; not microscopic global linearity.\\
Tunnelling/WKB & \statusConditional{} & C & Standard WKB in a separately supplied low-energy interacting-envelope reduction and barrier; a generic local potential does not preserve the free spectral branch automatically.\\
Hydrogen gross $1/n^2$ spectrum & \statusConditional{} & C & Free N2 kinetic bridge + $J_0$ + sourced Maxwell--Coulomb interaction + additional weakly bound two-body/branch reduction and $Q=1$ identification. $g_C$ is dimensionful, $\alpha_{\rm eff}=g_C/(J_0c)$; exact interacting branch selection and native electromagnetic current/source realization remain Open.\\
Decoherence & \statusModelDerived{} & C & Exact reduced-coherence identity in the supplied factorized system/environment model with its coupling and normalized environment states; no native coupling, pointer basis or outcome selection.\\
Zeno suppression & \statusConditional{} & C & Finite spectral variance yields short-interval survival expansion; ideal repeated rank-one survival projections further require a conditional state-reset/interrogation rule. Native apparatus/readout and anti-Zeno rate Open.\\
Framed/double-cover $4\pi$ return & \statusConditional{} & N2 & Mathematical/framed admissibility; ordinary single-valued tensor transport alone does not force the spinorial sign; physical matter-sector selection Open.\\
Physical spinorial-sector selector & \statusOpen{} & O & Not supplied by the one-body Clifford algebra.\\
Many-closure exchange sign & \statusSharpOpen{} & O & One-body $4\pi$ return does not select antisymmetry.\\
Pauli exclusion once antisymmetry supplied & \statusConditional{} & C & Exact downstream consequence once antisymmetry is supplied; the native exchange-sign/statistics bridge remains Open.\\
Literal local-steering origin of the complex unit & \statusRetracted{} & N2 & The tested identification $J_hv=v\times\widehat h$ does not select the shared quantum complex unit; independent phase/frame coupling candidates are not excluded.\\
Generic finite-wave action-current and normalization bridge & \statusSharpOpen{} & O & Unique positive physical current, modal provenance and persistence through nonlinear closure formation are not supplied by the resolved quadratic branch or its arbitrary scale $J_0$.\\
Physical complex structure / positive-frequency selection & \statusSharpOpen{} & O & Effective quadratures do not select the physical free-sector vacuum, shared complex unit or multi-system composition. The local steering-origin identification was tested and retracted.\\
Conjugate-closure / antimatter interpretation & \statusConjectural{} / \statusConditional{} & O & Candidate physical paired-branch structure; equal mass only if paired branches have equal invariant return energy. Algebraic conjugation does not derive physical $C$, opposite gauge charge/current, annihilation dynamics, CPT, or baryogenesis.\\
Born weights / apparatus / definite outcome & \statusOpen{} / \statusDeferred{} & O & Native preparation measure, detector coupling and context-dependent record-selection frequencies unproved; completion belongs to the separate publication-deferred measurement research programme.\\
Free Clifford stationary localization & \statusNegative{} & N2 & No nonzero $L^2$ stationary eigenstate in the stated free translation-invariant class.\\
Three families from Clifford/Weyl dimension & \statusNegative{} & N2 & Module dimensions are powers of two; family mechanism lies elsewhere.\\
\bottomrule
\end{longtable}
\endgroup

\newpage

\printbibliography[title={References}]


\end{document}
